\chapter{Introdução}
\label{cap:introducao}

% -----------------------------------------------------------------
% Objetivo do capítulo:
% Apresentar o contexto industrial, o problema dos rompimentos, a
% motivação para uma ferramenta de apoio à decisão e a estrutura
% geral do trabalho. A introdução deve responder, em ordem:
%   1. Qual é o cenário industrial?
%   2. Qual é o problema?
%   3. Por que ele é difícil?
%   4. O que a pesquisa propõe?
%   5. Quais são objetivos e contribuições?
%   6. Como o documento está organizado?
% -----------------------------------------------------------------


\section{Contexto industrial}
\label{sec:contexto_industrial}

% Base: docs/dissertacao/08_contexto_industrial.md
% Pontos a cobrir:
%   - Cadeia produtiva do fio esmaltado (trefilação, recozedor,
%     esmaltação, cestos, bobinador).
%   - Caráter contínuo do processo e o efeito de uma interrupção.
%   - Aplicações finais (motores, transformadores, bobinas).
%   - Multiplicidade de causas possíveis para um rompimento.


\section{Problema de pesquisa}
\label{sec:problema_pesquisa}

% Pergunta central:
% Como combinar sinais operacionais de máquina e evidências
% estatísticas históricas de rompimentos para apoiar a inferência
% da provável origem de eventos de rompimento em uma fábrica de
% fios?


\section{Proposição central}
\label{sec:proposicao_central}

% A integração entre um módulo de detecção de anomalias em
% variáveis operacionais e um módulo de análise estatística
% histórica de rompimentos por lote, material, fornecedor e
% padrão de cruzamento entre máquinas pode produzir um
% diagnóstico mais consistente, explicável e útil que abordagens
% isoladas baseadas em apenas um desses eixos.


\section{Objetivos}
\label{sec:objetivos}

\subsection{Objetivo geral}
\label{subsec:objetivo_geral}

% Desenvolver e avaliar uma ferramenta híbrida de apoio à
% decisão para inferir a provável origem de rompimentos em uma
% fábrica de fios, combinando detecção de anomalias operacionais
% com análise estatística histórica de falhas.

\subsection{Objetivos específicos}
\label{subsec:objetivos_especificos}

\begin{enumerate}[noitemsep,nosep,labelindent=\parindent,leftmargin=*,label={\alph*}) ]
	\item Estruturar uma base histórica de rompimentos e operação com qualidade suficiente para análise.
	\item Modelar o comportamento operacional esperado das máquinas e identificar desvios por meio de técnicas de detecção de anomalias.
	\item Construir indicadores estatísticos históricos de risco por lote, material, fornecedor e máquina, considerando exposição operacional e taxas relativas de rompimento.
	\item Definir uma lógica de integração entre evidências operacionais e históricas para produzir um diagnóstico da provável origem do rompimento.
	\item Avaliar a coerência técnica e a utilidade prática da ferramenta em dados reais de operação fabril.
	\item Produzir uma representação explicável do diagnóstico, de modo que o resultado possa ser interpretado por manutenção, qualidade e engenharia.
\end{enumerate}


\section{Justificativa}
\label{sec:justificativa}

% Pontos a cobrir:
%   - Relevância aplicada (perdas, retrabalho, dependência de
%     especialista).
%   - Lacuna técnica (sistemas que integram ML e estatística
%     histórica com interpretabilidade em chão de fábrica).
%   - Relevância científica (arquitetura híbrida e explicável
%     aplicada a domínio industrial real).


\section{Contribuições esperadas}
\label{sec:contribuicoes}

\begin{enumerate}[noitemsep,nosep,labelindent=\parindent,leftmargin=*,label={\alph*}) ]
	\item Proposição de uma arquitetura híbrida para análise de rompimentos em ambiente industrial real.
	\item Integração entre aprendizado de máquina não supervisionado e estatística histórica aplicada ao diagnóstico operacional.
	\item Mecanismo explicável de apoio à decisão para diferenciar indícios de falha de equipamento e suspeita de problema associado a matéria-prima.
	\item Estruturação de indicadores históricos operacionais reutilizáveis em rotinas de monitoramento e melhoria contínua.
	\item Geração de conhecimento aplicado para o contexto industrial da fabricação de fios, com potencial de adaptação para outros processos com falhas multifatoriais.
\end{enumerate}


\section{Delimitação metodológica}
\label{sec:delimitacao}

% O trabalho deve ser apresentado como sistema de inferência da
% provável origem do rompimento, não como prova causal
% definitiva. Essa delimitação alinha a promessa acadêmica ao
% que o método entrega: apoio quantitativo e explicável à
% decisão em ambiente de incerteza.


\section{Organização do trabalho}
\label{sec:organizacao}

% Sugestão de parágrafo final, a redigir depois que todos os
% capítulos estiverem com título estável:
%   - Capítulo 2 apresenta a fundamentação teórica.
%   - Capítulo 3 descreve a metodologia adotada.
%   - Capítulo 4 detalha o desenvolvimento da ferramenta.
%   - Capítulo 5 traz os resultados e a discussão.
%   - Capítulo 6 apresenta as conclusões e os trabalhos futuros.
